% 2.7 — brief en documento/investigacion/2-7-stack-desarrollo.md
\subsection{TECNOLOGÍAS DEL STACK DE DESARROLLO}

Esta sección presenta las tecnologías con las que se construye el sistema y justifica su elección. Conviene distinguir dos grupos. El stack del ERP (NestJS, Prisma, PostgreSQL y Angular) está fijado por el sistema existente y este trabajo lo toma como dado. El stack del módulo predictivo (Python, DuckDB, Polars, Scikit-learn, XGBoost y FastAPI) es una decisión de diseño de este proyecto, todavía no implementada, y es la que requiere justificación. Keycloak y el servidor MCP abarcan ambos grupos, y Power BI solo el ERP; los tres se tratan en las secciones 2.8 a 2.10.

\subsubsection{PYTHON Y SU ECOSISTEMA ML}

Python es un lenguaje de programación interpretado, interactivo y orientado a objetos, que admite además los paradigmas procedural y funcional, es portable entre Linux, macOS y Windows, y es extensible en C o C++ \parencite{PythonFAQ}. Esta última propiedad explica la forma que tomó su ecosistema científico: las bibliotecas numéricas exponen una interfaz en Python sobre un núcleo compilado. La base de ese ecosistema es NumPy, que \textcite{Harris2020} presentan como la biblioteca principal de programación con arreglos del lenguaje y como el fundamento sobre el que se construye el resto de las bibliotecas científicas, con el papel adicional de capa de interoperabilidad entre ellas. La programación con arreglos ofrece una sintaxis compacta para operar sobre vectores y matrices completos en lugar de elemento por elemento; ese estilo vectorizado reaparece, con otra implementación, en los motores de DuckDB y Polars.

Sobre NumPy se apoya pandas, presentada por \textcite{McKinney2010} como una biblioteca de estructuras de datos para el cómputo estadístico, motivada por los problemas prácticos de trabajar con los conjuntos de datos propios de las finanzas y la estadística. Sus dos estructuras primarias son la \texttt{Series}, unidimensional y etiquetada, y el \texttt{DataFrame}, tabular y con columnas de tipos heterogéneos; la biblioteca resuelve de forma nativa el manejo de valores faltantes, la alineación por etiquetas, las operaciones de agrupamiento y combinación, y funciones específicas para series temporales como ventanas móviles, desplazamientos y rezagos \parencite{PandasOverview}. En este proyecto pandas cumple dos funciones: la exploración del histórico en el Capítulo III y el puente hacia las bibliotecas de modelado, que aceptan \texttt{DataFrame} como entrada.

La capa de modelado la forman Scikit-learn y XGBoost. \textcite{Pedregosa2011} presentan Scikit-learn como una biblioteca de aprendizaje automático que pone el énfasis en la facilidad de uso, el desempeño, la documentación y la consistencia de su interfaz, con licencia BSD y dependencias mínimas. XGBoost, cuyo algoritmo se describió en la sección 2.2, se emplea aquí como biblioteca: su paquete para Python incluye estimadores compatibles con la interfaz de Scikit-learn (\texttt{XGBClassifier} y \texttt{XGBRegressor}), lo que permite usarlos dentro de los \textit{pipelines} y de las rutinas de validación cruzada de esa biblioteca \parencite{Chen2016, XGBoostDocs}. La elección de Python para el módulo predictivo responde a que toda la cadena (arreglos, tablas, estimadores y servicio de predicción) se resuelve en un mismo lenguaje, y a que los tres componentes restantes de esa cadena, DuckDB, Polars y FastAPI, se usan desde Python. La versión concreta del intérprete y de cada biblioteca se fija en el Capítulo IV al construir el pipeline.

\subsubsection{DUCKDB Y POLARS}

La decisión de no ejecutar el análisis sobre la base de datos de producción del ERP tiene un fundamento de arquitectura. \textcite{Stonebraker2005} sostienen que la arquitectura tradicional de los sistemas de gestión de bases de datos, diseñada y optimizada para el procesamiento de datos de negocio, se ha usado para sostener aplicaciones con características y requisitos muy distintos, y argumentan que la idea de un motor único para todos los usos ya no es aplicable: el mercado se fragmenta en motores independientes, y el almacén de datos es uno de los dos ejemplos con que sustentan el argumento. Trasladado a este proyecto, el procesamiento transaccional en línea (\textit{online transaction processing}, OLTP) del ERP y el procesamiento analítico en línea (\textit{online analytical processing}, OLAP) del pipeline Medallion tienen requisitos distintos y se resuelven con motores distintos. El argumento es de arquitectura, no de rendimiento medido en este trabajo.

El motor analítico elegido es DuckDB. \textcite{Raasveldt2019} lo presentan como un sistema de gestión de datos diseñado para ejecutar consultas SQL analíticas embebido en otro proceso, una necesidad que la popularidad de SQLite hacía evidente y que ningún sistema cubría para cargas analíticas; se distribuye como software libre con licencia permisiva. Según su documentación, no tiene dependencias externas, usa un motor de ejecución columnar y vectorizado, que procesa un lote de valores por operación, y es extensible en tipos, funciones y formatos \parencite{DuckDBWhy}. Dos extensiones justifican su papel en la capa Silver: la de MySQL permite consultar directamente una instancia MySQL en ejecución y cargar sus tablas en DuckDB, y la de PostgreSQL hace lo propio con esa base, ambas mediante la sentencia \texttt{ATTACH} \parencite{DuckDBMySQL, DuckDBPostgres}. Un solo motor SQL, embebido en el proceso Python del ETL y sin servidor adicional que desplegar, lee así el histórico del portal anterior, los proyectos activos del ERP nuevo y los archivos de la capa Bronze. La documentación de la extensión MySQL no indica versiones mínimas soportadas, por lo que su compatibilidad con MySQL 5.7, la versión del portal anterior, se verifica empíricamente en el Capítulo IV\@.

Para la transformación de la capa Gold se prevé utilizar Polars, una biblioteca de \texttt{DataFrame} con núcleo escrito en Rust y disponible para Python \parencite{PolarsDocs}. Su documentación declara ejecución paralela entre los núcleos del procesador, un motor de consultas vectorizado, una interfaz de flujo (\textit{streaming}) para procesar resultados sin tener todos los datos en memoria, y soporte de Apache Arrow. Su interfaz perezosa (\textit{lazy}) difiere la ejecución hasta que se solicita el resultado, lo que le permite filtrar filas y seleccionar columnas lo antes posible durante la lectura (empuje de predicados y de proyecciones); la documentación recomienda preferirla salvo en trabajo exploratorio \parencite{PolarsLazy}. Esas propiedades, y no una comparación de velocidad con pandas que no se ha verificado con fuentes independientes, motivan la elección. El traspaso entre las dos capas es directo: DuckDB consulta de forma nativa los \texttt{DataFrame} de Polars presentes en el ámbito de ejecución y devuelve sus resultados como \texttt{DataFrame} de Polars mediante la integración con Arrow, sin serializar a disco \parencite{DuckDBPolars}.

\subsubsection{FASTAPI COMO MICROSERVICIO}

El servicio de predicción se construye sobre el estándar ASGI (\textit{Asynchronous Server Gateway Interface}), que su especificación define como sucesor de WSGI, el estándar de compatibilidad entre servidores web, marcos de trabajo y aplicaciones en Python. Una aplicación WSGI es un único invocable síncrono que recibe una petición y devuelve una respuesta, lo que no admite conexiones de larga duración; ASGI extiende ese contrato al Python asíncrono \parencite{ASGISpec}. Sobre el estándar se apoya Starlette, un marco ligero para construir servicios web asíncronos, con pocas dependencias obligatorias y código completamente anotado con tipos \parencite{StarletteDocs}, y lo ejecuta un servidor ASGI como Uvicorn \parencite{FastAPIASGI}.

La segunda pieza es Pydantic, biblioteca de validación de datos en la que el esquema de validación y serialización se controla mediante anotaciones de tipo; su lógica central está escrita en Rust, admite un modo estricto sin conversión de tipos y uno laxo con coerción, y emite JSON Schema para integrarse con otras herramientas \parencite{PydanticDocs}. FastAPI combina ambas: según su documentación, se apoya en Starlette para la parte web y en Pydantic para la parte de datos, y está basado en los estándares abiertos OpenAPI y JSON Schema, lo que le permite generar documentación interactiva automática (Swagger UI y ReDoc) y habilitar la generación de clientes en otros lenguajes \parencite{FastAPIDocs}. FastAPI es, en consecuencia, un marco ASGI \parencite{FastAPIASGI}.

Tres razones justifican su elección para el módulo predictivo. Primero, los modelos entrenados con Scikit-learn y XGBoost viven en Python; servirlos desde Python evita reimplementarlos o trasladarlos a otro entorno de ejecución. Segundo, Pydantic define el contrato de entrada y salida del endpoint de predicción (tipos, rangos y valores faltantes admitidos) y FastAPI lo publica como especificación OpenAPI, que es el contrato que el backend NestJS consume para tipar su cliente. Tercero, el modelo asíncrono de ASGI permite que la obtención de las claves públicas de Keycloak para validar tokens (sección 2.8) y otras operaciones de entrada y salida no bloqueen al proceso mientras espera la respuesta. El servicio se despliega como un contenedor independiente del ERP, de modo que el modelo puede actualizarse sin redesplegar el backend. Las cifras de rendimiento y productividad que la documentación de FastAPI publica sobre sí misma no se toman como evidencia, por tratarse de autoevaluaciones sin estudio independiente.

\subsubsection{NESTJS Y ANGULAR}

El backend del ERP usa NestJS, un marco para construir aplicaciones de servidor en Node.js que está escrito en TypeScript y lo soporta por completo, combina elementos de la programación orientada a objetos, funcional y reactiva funcional, y utiliza Express como servidor HTTP por defecto \parencite{NestJSDocs}; el ERP corre NestJS 11 sobre Express y Node 22. Su unidad de organización es el módulo, una clase anotada con el decorador \texttt{@Module()}: toda aplicación tiene al menos un módulo raíz desde el que Nest construye el grafo de la aplicación, y cada módulo encapsula sus proveedores, de modo que solo pueden inyectarse los propios o los exportados explícitamente por módulos importados \parencite{NestJSModules}. Esa encapsulación permite agregar al ERP, como módulos separados, la integración con el servicio de predicción y la API para Power BI (sección 2.10); el servidor MCP, en cambio, corre como proceso independiente (sección 2.9).

El acceso a datos lo resuelve Prisma, un ORM (\textit{object-relational mapping}) para Node.js y TypeScript que ofrece acceso a la base con seguridad de tipos, migraciones y un editor visual de datos; su esquema es la única fuente de verdad de los modelos de base de datos y de aplicación, y las consultas se validan en tiempo de compilación \parencite{PrismaDocs7}. El ERP usa Prisma 7.8, versión en la que el adaptador de controlador es obligatorio para las conexiones directas y la configuración vive en \texttt{prisma.config.ts}. La base es PostgreSQL 17, un sistema de gestión de bases de datos objeto-relacional que soporta gran parte del estándar SQL y ofrece integridad transaccional y control de concurrencia multiversión \parencite{PostgreSQLDocs}; esas son las propiedades OLTP que, según el argumento de \textcite{Stonebraker2005} recogido en la subsección anterior, conviene no mezclar con la carga analítica.

El frontend es Angular 21.2, un marco web mantenido por un equipo de Google que organiza el código en componentes encapsulados e incorpora un modelo de reactividad de grano fino basado en \textit{signals}, inyección de dependencias, enrutamiento y formularios con validación \parencite{AngularOverview}. El ERP lo usa con renderizado del lado del servidor (\textit{server-side rendering}, SSR), que genera en el servidor una página HTML completa para cada petición y luego la reutiliza en el cliente mediante hidratación \parencite{AngularSSR}. Para una aplicación interna con inicio de sesión, el beneficio relevante no es el posicionamiento en buscadores sino el primer render rápido del dashboard donde se muestran las predicciones. La interfaz se compone con PrimeNG 21.1 y Tailwind 4.1. Todos los componentes se despliegan con Docker, que empaqueta y ejecuta cada aplicación en un contenedor aislado que sirve como unidad de distribución \parencite{DockerDocs}; el monorepo define los servicios en un archivo de Docker Compose, y la portabilidad de las imágenes permite prever la migración desde el servidor propio de la empresa a la nube como trabajo futuro. El cuadro~\ref{tab:stack-desarrollo} resume el stack completo.

\begin{table}[H]
  \centering
  \caption{Tecnologías del sistema por capa de la arquitectura}\label{tab:stack-desarrollo}
  \begin{tblr}{colspec={X[1.04,l] X[0.96,l] X[1.39,l] X[0.61,l]}}
    \textbf{Tecnología} & \textbf{Capa} & \textbf{Rol} & \textbf{Versión} \\
    \SetCell[c=4]{l} \textit{ERP existente} & & & \\
    NestJS (Express, Node) & Backend del ERP & API y módulos de negocio; integración con FastAPI y API para Power BI; consumido por el servidor MCP & 11 (Node 22) \\
    Prisma & Backend del ERP & ORM y migraciones & 7.8 \\
    PostgreSQL & Base OLTP del ERP & Datos transaccionales; fuente de inferencia & 17 \\
    Angular (SSR) & Frontend del ERP & Interfaz de usuario y dashboard de predicciones & 21.2 \\
    PrimeNG / Tailwind & Frontend del ERP & Componentes de interfaz y estilos & 21.1 / 4.1 \\
    PHP + MySQL & Portal anterior & Fuente del histórico (entrenamiento) & MySQL 5.7 \\
    \SetCell[c=4]{l} \textit{Módulo predictivo (por implementar)} & & & \\
    Python (NumPy, pandas) & Pipeline y modelos & Lenguaje del ETL, exploración y modelado & Cap.~IV \\
    DuckDB & Silver & Motor OLAP embebido; lee MySQL, PostgreSQL y archivos & Cap.~IV \\
    Polars & Gold & Ingeniería de características & Cap.~IV \\
    Scikit-learn / XGBoost & Modelos & Pipelines, validación cruzada y \textit{gradient boosting} & Cap.~IV \\
    FastAPI (Starlette, Pydantic, Uvicorn) & Servicio de predicción & Endpoint de inferencia con contrato OpenAPI & Cap.~IV \\
    \SetCell[c=4]{l} \textit{Transversales} & & & \\
    Keycloak & Identidad & SSO con OAuth 2.0 / OpenID Connect (sección 2.8) & Cap.~IV \\
    Servidor MCP & Integración con IA & Tools sobre el ERP y las predicciones (sección 2.9) & Cap.~IV \\
    Power BI & Inteligencia de negocios & Reportes descriptivos vía API NestJS (sección 2.10) & --- \\
    Docker / Docker Compose & Despliegue & Contenedores en servidor propio; migración a nube como trabajo futuro & --- \\
  \end{tblr}
  \fuente{Elaboración propia, 2026, en base a la documentación y los repositorios del ERP de ISI Mustang.}
\end{table}
